%-----------------------------------------------------------------------------------------
% Instruções para formação
% Crie um arquivo CSV chamado 'formacao.csv' com as colunas:
% TituloAtividade;Periodo;InstituicaoEnsino
%
% Exemplo:
% Doutorado em Ciências Biológicas, área de Zoologia;2014 a 2018;Universidade Exemplo
%
% Salve no mesmo diretório deste documento.
% Dica: abra o formacao.csv no Excel, edite e salve como "CSV UTF-8" (veja o README).
%-----------------------------------------------------------------------------------------

\pgfplotstabletypeset[
    col sep=semicolon,
    header=true,        % a primeira linha do CSV é o cabeçalho
    string type,
    font=\normalsize,
    every head row/.style={before row=\toprule, after row=\midrule},
    every last row/.style={after row=\bottomrule},
    every even row/.style={before row={\rowcolor[gray]{0.85}}},
    every odd row/.style={before row={\rowcolor[gray]{0.95}}},
    columns={[index]0,[index]1,[index]2},
    display columns/0/.style={
        column name={\textbf{Título ou atividade}},
        column type={>{\raggedright\arraybackslash}m{4.5cm}}
    },
    display columns/1/.style={
        column name={\textbf{Início--Término}},
        column type={>{\raggedright\arraybackslash}m{3.0cm}}
    },
    display columns/2/.style={
        column name={\textbf{Instituição de ensino}},
        column type={>{\raggedright\arraybackslash}m{7.1cm}}
    }
]{formacao.csv}

\vspace{1em}

\emph{\large 1.1) Formação -- Informações Adicionais}

\begin{itemize}
    \item \textbf{AAAA--AAAA:} Período dedicado a atividade relevante que não aparece na tabela acima (por exemplo, preparação para concursos públicos, estágio fora da academia).
    \item \textbf{Mês a mês de AAAA:} Afastamento ou licença que explique lacuna no currículo (por exemplo, licença-maternidade/paternidade, licença médica).
\end{itemize}
